@@M0@@ to @@M1@@; denote by @@M2@@ the homomorphism defined by the sequence @@M3@@, so that @@M4@@ is deduced from @@M5@@ by base change from @@M6@@ to @@M7@@. Since the sequence @@M8@@ is regular (in the classical sense), the homomorphism @@M9@@ is regular, and the hypotheses of the lemma hold for the base change @@M10@@. Now the @@M11@@ are free of finite type over @@M12@@, so we can apply the results of a). Consequently @@M13@@, and @@M14@@ is free of finite type.

\sgaStableTarget{sga6.e07.prop.1-3}\textbf{Proposition 1.3.} \uline{Let} @@M15@@ \uline{be a ring}, @@M16@@ \uline{a projective} @@M17@@-\uline{module of finite type}, @@M18@@ \uline{a regular homomorphism}, @@M19@@. \uline{Then the ideal} @@M20@@ \uline{is quasi-regular} (EGA IV 16.9.7), \uline{that is} (EGA IV 16.9.8):
\begin{enumerate}[label=\roman*)]
\item @@M21@@ \uline{is of finite type}.
\item @@M22@@ \uline{is a projective} @@M23@@-\uline{module of finite type};
\item \uline{the canonical homomorphism}
\end{enumerate}
\sgaStableTarget{sga6.e07.eq.1-3-1}@@M24@@
\uline{is an isomorphism}.

We saw in the proof of 1.2 b) that @@M25@@ is projective of finite type. Resume the notation of 1.2 b) (since the question is local on @@M26@@, we again suppose @@M27@@ free of finite type). Then the canonical homomorphism (1.3.1) is deduced from the analogous homomorphism for @@M28@@ by the base change @@M29@@, by the relations 1.2. The latter is an isomorphism, so the same is true of (1.3.1).

\sgaStableTarget{sga6.e07.def.1-4}\textbf{Definition 1.4.} \uline{Let} @@M30@@ \uline{be a commutatively locally ringed site}, @@M31@@ \uline{an} @@M32@@-\uline{Module locally free of finite type}, @@M33@@ \uline{an} @@M34@@-\uline{linear homomorphism}. \uline{We say that} @@M35@@ \uline{is} regular \uline{if the Koszul complex} @@M36@@ \uline{defined by} @@M37@@ \uline{is acyclic in degrees} @@M38@@. \uline{Let} @@M39@@ \uline{be an Ideal of} @@M40@@; \uline{we say that} @@M41@@ \uline{is}
